% Options for packages loaded elsewhere
\PassOptionsToPackage{unicode}{hyperref}
\PassOptionsToPackage{hyphens}{url}
\PassOptionsToPackage{dvipsnames,svgnames,x11names}{xcolor}
%
\documentclass[
  10pt,
]{article}
\usepackage{amsmath,amssymb}
\usepackage{iftex}
\ifPDFTeX
  \usepackage[T1]{fontenc}
  \usepackage[utf8]{inputenc}
  \usepackage{textcomp} % provide euro and other symbols
\else % if luatex or xetex
  \usepackage{unicode-math} % this also loads fontspec
  \defaultfontfeatures{Scale=MatchLowercase}
  \defaultfontfeatures[\rmfamily]{Ligatures=TeX,Scale=1}
\fi
\usepackage{lmodern}
\ifPDFTeX\else
  % xetex/luatex font selection
\fi
% Use upquote if available, for straight quotes in verbatim environments
\IfFileExists{upquote.sty}{\usepackage{upquote}}{}
\IfFileExists{microtype.sty}{% use microtype if available
  \usepackage[]{microtype}
  \UseMicrotypeSet[protrusion]{basicmath} % disable protrusion for tt fonts
}{}
\makeatletter
\@ifundefined{KOMAClassName}{% if non-KOMA class
  \IfFileExists{parskip.sty}{%
    \usepackage{parskip}
  }{% else
    \setlength{\parindent}{0pt}
    \setlength{\parskip}{6pt plus 2pt minus 1pt}}
}{% if KOMA class
  \KOMAoptions{parskip=half}}
\makeatother
\usepackage{xcolor}
\usepackage[margin=1in]{geometry}
\usepackage{color}
\usepackage{fancyvrb}
\newcommand{\VerbBar}{|}
\newcommand{\VERB}{\Verb[commandchars=\\\{\}]}
\DefineVerbatimEnvironment{Highlighting}{Verbatim}{commandchars=\\\{\}}
% Add ',fontsize=\small' for more characters per line
\newenvironment{Shaded}{}{}
\newcommand{\AlertTok}[1]{\textcolor[rgb]{1.00,0.00,0.00}{\textbf{#1}}}
\newcommand{\AnnotationTok}[1]{\textcolor[rgb]{0.38,0.63,0.69}{\textbf{\textit{#1}}}}
\newcommand{\AttributeTok}[1]{\textcolor[rgb]{0.49,0.56,0.16}{#1}}
\newcommand{\BaseNTok}[1]{\textcolor[rgb]{0.25,0.63,0.44}{#1}}
\newcommand{\BuiltInTok}[1]{\textcolor[rgb]{0.00,0.50,0.00}{#1}}
\newcommand{\CharTok}[1]{\textcolor[rgb]{0.25,0.44,0.63}{#1}}
\newcommand{\CommentTok}[1]{\textcolor[rgb]{0.38,0.63,0.69}{\textit{#1}}}
\newcommand{\CommentVarTok}[1]{\textcolor[rgb]{0.38,0.63,0.69}{\textbf{\textit{#1}}}}
\newcommand{\ConstantTok}[1]{\textcolor[rgb]{0.53,0.00,0.00}{#1}}
\newcommand{\ControlFlowTok}[1]{\textcolor[rgb]{0.00,0.44,0.13}{\textbf{#1}}}
\newcommand{\DataTypeTok}[1]{\textcolor[rgb]{0.56,0.13,0.00}{#1}}
\newcommand{\DecValTok}[1]{\textcolor[rgb]{0.25,0.63,0.44}{#1}}
\newcommand{\DocumentationTok}[1]{\textcolor[rgb]{0.73,0.13,0.13}{\textit{#1}}}
\newcommand{\ErrorTok}[1]{\textcolor[rgb]{1.00,0.00,0.00}{\textbf{#1}}}
\newcommand{\ExtensionTok}[1]{#1}
\newcommand{\FloatTok}[1]{\textcolor[rgb]{0.25,0.63,0.44}{#1}}
\newcommand{\FunctionTok}[1]{\textcolor[rgb]{0.02,0.16,0.49}{#1}}
\newcommand{\ImportTok}[1]{\textcolor[rgb]{0.00,0.50,0.00}{\textbf{#1}}}
\newcommand{\InformationTok}[1]{\textcolor[rgb]{0.38,0.63,0.69}{\textbf{\textit{#1}}}}
\newcommand{\KeywordTok}[1]{\textcolor[rgb]{0.00,0.44,0.13}{\textbf{#1}}}
\newcommand{\NormalTok}[1]{#1}
\newcommand{\OperatorTok}[1]{\textcolor[rgb]{0.40,0.40,0.40}{#1}}
\newcommand{\OtherTok}[1]{\textcolor[rgb]{0.00,0.44,0.13}{#1}}
\newcommand{\PreprocessorTok}[1]{\textcolor[rgb]{0.74,0.48,0.00}{#1}}
\newcommand{\RegionMarkerTok}[1]{#1}
\newcommand{\SpecialCharTok}[1]{\textcolor[rgb]{0.25,0.44,0.63}{#1}}
\newcommand{\SpecialStringTok}[1]{\textcolor[rgb]{0.73,0.40,0.53}{#1}}
\newcommand{\StringTok}[1]{\textcolor[rgb]{0.25,0.44,0.63}{#1}}
\newcommand{\VariableTok}[1]{\textcolor[rgb]{0.10,0.09,0.49}{#1}}
\newcommand{\VerbatimStringTok}[1]{\textcolor[rgb]{0.25,0.44,0.63}{#1}}
\newcommand{\WarningTok}[1]{\textcolor[rgb]{0.38,0.63,0.69}{\textbf{\textit{#1}}}}
\setlength{\emergencystretch}{3em} % prevent overfull lines
\providecommand{\tightlist}{%
  \setlength{\itemsep}{0pt}\setlength{\parskip}{0pt}}
\setcounter{secnumdepth}{-\maxdimen} % remove section numbering
\ifLuaTeX
  \usepackage{selnolig}  % disable illegal ligatures
\fi
\usepackage{bookmark}
\IfFileExists{xurl.sty}{\usepackage{xurl}}{} % add URL line breaks if available
\urlstyle{same}
\hypersetup{
  colorlinks=true,
  linkcolor={Maroon},
  filecolor={Maroon},
  citecolor={Blue},
  urlcolor={blue},
  pdfcreator={LaTeX via pandoc}}

\author{}
\date{}

\begin{document}

\section{DDTA minimal BAE ontology candidate -
R2}\label{ddta-minimal-bae-ontology-candidate---r2}

\textbf{Status:} STRONG CANDIDATE / NOT ACCEPTED / BA1 OPEN

\textbf{Derived by:} BA1-T1 + BA1-T2

\textbf{Repository baseline reviewed:}
\texttt{f05fbb7b253392e158a1062df2614b177c13d43e}

\textbf{Supersedes for active BA1 reasoning:}
\texttt{BA1\_MINIMAL\_BAE\_ONTOLOGY\_CANDIDATE\_R1.md}

R1 remains immutable research history.

\subsection{Candidate ontology}\label{candidate-ontology}

The current minimal candidate contains exactly two semantic identity
families:

\begin{Shaded}
\begin{Highlighting}[]
\NormalTok{BAReferent}
\NormalTok{BAProposition}
\end{Highlighting}
\end{Shaded}

No domain-specific subtype split is accepted yet.

\subsection{\texorpdfstring{\texttt{BAReferent}}{BAReferent}}\label{bareferent}

A \texttt{BAReferent} is an independently identifiable unit of
methodology-neutral shared project meaning that must be reusable across
one or more analytical propositions, projections or governed baselines.

A referent may denote structural or human/external meaning, information
or resources, behavior or events, contracts, storage, boundaries, states
or contexts, or other project meaning.

Those semantic kinds may be important classifications, but BA1-T2 does
not promote them to dedicated first-class metaclasses.

When a behavior, event, milestone, state or other meaning must be
independently reused or constrained, that meaning can itself receive
\texttt{BAReferent} identity.

\subsection{\texorpdfstring{\texttt{BAProposition}}{BAProposition}}\label{baproposition}

A \texttt{BAProposition} is an independently identifiable
methodology-neutral analytical assertion about one or more
\texttt{BAReferent} identities, with enough identity to support
origin/provenance, diagnosis, reuse and change disposition.

A proposition may express project meaning such as:

\begin{itemize}
\tightlist
\item
  relation/dependency;
\item
  behavior/action/transition semantics;
\item
  responsibility/externality;
\item
  correlation;
\item
  state/condition/result semantics;
\item
  applicability;
\item
  constraint/property;
\item
  lifecycle/failure/concurrency semantics.
\end{itemize}

The exact proposition structure, predicate/action vocabulary, role
vocabulary and cardinalities remain BA2 questions.

\subsection{Refined family boundary}\label{refined-family-boundary}

A \texttt{BAProposition} is not itself the project meaning that it
describes.

If project meaning must be independently reused or qualified, that
meaning receives \texttt{BAReferent} identity and propositions state its
semantics.

Therefore proposition-as-project-target is \textbf{not required by the
minimal candidate}.

Proposition identity still exists so that BA3 can attach
origin/provenance, diagnostic state and change disposition to the
assertion itself without confusing those concerns with referent
identity.

\subsection{Minimal picture}\label{minimal-picture}

\begin{Shaded}
\begin{Highlighting}[]
\NormalTok{Base Analysis for one governed baseline}

\NormalTok{BAReferent *}
\NormalTok{    \^{}}
\NormalTok{    | referenced by}
\NormalTok{    |}
\NormalTok{BAProposition *}

\NormalTok{project{-}semantic reuse/qualification {-}\textgreater{} BAReferent}
\NormalTok{assertion provenance/diagnostic state {-}\textgreater{} BAProposition metadata in BA3}
\end{Highlighting}
\end{Shaded}

\texttt{BAE} remains an umbrella term, not an additional required
metaclass.

\subsection{Split-or-collapse result}\label{split-or-collapse-result}

BA1-T2 did not force dedicated first-class roots for:

\begin{itemize}
\tightlist
\item
  Participant / Actor;
\item
  Component / Capability;
\item
  Information / Resource;
\item
  Behavior / Event / Transition;
\item
  Interface / Connection / Channel;
\item
  Store / Persistence;
\item
  Contract;
\item
  Boundary / Externality;
\item
  State / Mode / Context;
\item
  Dependency / Allocation;
\item
  Property / Constraint;
\item
  DataFlow / Interaction.
\end{itemize}

A concrete instance of any of those meanings may still be a
\texttt{BAReferent} when independent identity is needed.

The opposite collapse into one undifferentiated semantic family is
rejected because referent identity and assertion
identity/origin/evolution are independently required.

A single programming class, table or serialization container with a
mandatory semantic discriminator would be only a representation choice
and would not collapse the two semantic families.

\subsection{Document-layer boundary}\label{document-layer-boundary}

Document types remain source-layer governed concepts and are not copied
into Base Analysis as BAE types:

\begin{itemize}
\tightlist
\item
  \texttt{MacroRequirement};
\item
  \texttt{Decision};
\item
  \texttt{Requirement};
\item
  \texttt{FunctionalRequirement};
\item
  \texttt{SpecializedRequirement};
\item
  \texttt{SecurityRequirement};
\item
  \texttt{NormativeClause}.
\end{itemize}

Governed documentation remains project authority.

\subsection{Deferred rather than
promoted}\label{deferred-rather-than-promoted}

The following remain required later responsibilities but are not domain
BAE metaclasses in this candidate:

\begin{itemize}
\tightlist
\item
  source locator and provenance materialization -\textgreater{} BA3;
\item
  grounded/derived/diagnostic materialization -\textgreater{} BA3;
\item
  identity/equivalence/lifecycle mechanics -\textgreater{} BA3;
\item
  views/projections -\textgreater{} BA4;
\item
  controlled vocabulary and lexical assistance -\textgreater{} BA5;
\item
  Base Analysis regression/closure -\textgreater{} BA6;
\item
  AnalysisRecord, Finding and method-specific outputs -\textgreater{}
  later analysis envelope.
\end{itemize}

\subsection{First-class split
criterion}\label{first-class-split-criterion}

A new first-class BAE type/family is justified only if concrete evidence
shows both:

\begin{enumerate}
\def\labelenumi{\arabic{enumi}.}
\tightlist
\item
  independent semantic identity is required across
  propositions/projections/change; and
\item
  reusable subtype-specific invariants cannot be represented honestly as
  classifications, roles, values or propositions over the existing
  families.
\end{enumerate}

Recurring usefulness, diagram convenience, implementation classes and
method-specific categories are insufficient.

\subsection{Falsification rule}\label{falsification-rule}

The two-family candidate fails if a concrete governed corpus or
method-neutral consumer requires a deferred semantic kind to have
independent subtype-specific identity/invariants that cannot be
represented honestly with \texttt{BAReferent\ +\ BAProposition}, or if
proposition identity can be removed without losing claim-level
provenance/change/diagnostic semantics.

\subsection{Next step}\label{next-step}

\texttt{BA1-T3\ -\ minimal\ BAE\ ontology\ closure\ review}.

BA1 remains open. BA2 remains not started.

\end{document}
